\documentclass[conference]{IEEEtran}
\IEEEoverridecommandlockouts
\usepackage{cite}
\usepackage{amsmath,amssymb,amsfonts}
\usepackage{algorithmic}
\usepackage{graphicx}
\usepackage{textcomp}
\usepackage{xcolor}
\def\BibTeX{{\rm B\kern-.05em{\sc i\kern-.025em b}\kern-.08em
    T\kern-.1667em\lower.7ex\hbox{E}\kern-.125emX}}

\begin{document}

\title{Project Checkpoint: Evaluating Audio-Visual Verification Robustness Against Deepfakes}

\author{\IEEEauthorblockN{Garvit Sharma}
\IEEEauthorblockA{\textit{Department of Computer Engineering} \\
\textit{San Jose State University}\\
San Jose, CA \\
garvit.sharma@sjsu.edu}
}

\maketitle

\begin{abstract}
This report outlines the initial datasets, exploratory data analysis, and baseline evaluation framework for assessing the vulnerability of the zero-shot audio-visual speaker verification architecture, specifically the Recursive Joint Cross-Attention (RJCA) model, against modern deepfake presentation attacks. We highlight a 100\% Attack Success Rate against sophisticated synthetic modalities during preliminary testing, quantify the theoretical weakness utilizing adversarial mathematics, and propose future defense frameworks for Phase 2.
\end{abstract}

\section{Refined Problem Statement and Motivation}
As biometric authentication systems shift towards multi-modal architectures (combining facial recognition with voice verification) to increase security, attackers are similarly pivoting towards highly synchronized deepfakes. Our core problem revolves around quantifying the security failures of these state-of-the-art authentication frameworks. 

We hypothesize that multi-modal systems relying intensely on lip-voice temporal synchronization (such as those powered by cross-attention mechanisms) are fundamentally critically vulnerable if they have not been actively trained with an adversarial discriminator block. This problem is exceptionally significant today because modern generative models (like Wav2Lip) can create synthetic identities with perfect audio-to-lip alignment, trivializing verification strategies relying solely on spatial-temporal consistency. 

\textcolor{blue}{Change in Scope: While our original proposal suggested developing a completely novel authentication network from scratch initially, our execution shifted slightly to first solidly establishing a vulnerable zero-shot baseline using the pre-existing published RJCA (Recursive Joint Cross-Attention) model. This ensures we have a strictly verified, mathematically calculated threshold constraint to measure any future spoofing-detection layers against explicitly.}

\section{Related Work and Literature Review}
In recent years, the rapid advancement of Generative Adversarial Networks (GANs) has given rise to sophisticated deepfake rendering networks such as \textbf{Wav2Lip} and \textbf{FSGAN}. While traditional uni-modal facial recognition systems trivially fail against these attacks, modern biometric security literature heavily advocates for multi-modal architectures to cross-reference integrity. Frameworks like the Recursive Joint Cross-Attention (RJCA) model leverage complex cross-modal attention mechanisms to dynamically weight visual lip movements against synchronized acoustic vocalizations in real-time. 

However, current verification literature largely operates under the naive assumption that if an attacker generates a visually pristine deepfake, the inherent sub-pixel mismatch between the acoustic waveform and the synthetic temporal lip rendering will automatically flag the verification logic. This report rigorously tests this hypothesis empirically against zero-shot execution environments.

\section{Mathematical Architecture of the RJCA Baseline}
To strictly formalize the vulnerability assessment, we established an inference environment utilizing a state-of-the-art \textbf{IResNet-18} deep residual network backbone for frame-wise spatial feature extraction, paired concurrently with an \textbf{ECAPA-TDNN} acoustic model for high-density vocal feature representations.

\subsection{Cross-Modal Attention Formulation}
The core fusion vector determining biometric identity relies on a joint cross-attention mechanism bridging the visual modality $V$ and audio modality $A$. By mapping queries $Q$, keys $K$, and values $V$ across the independent data domains, the underlying algorithm determines the physical synchronization overlap:
\begin{equation}
    \text{Attention}(Q, K, V) = \text{softmax}\left(\frac{QK^T}{\sqrt{d_k}}\right)V
\end{equation}
The resulting temporally-aligned features are ultimately mapped through a 1408-dimensional fully connected linear layer.

\subsection{Attentive Statistics Pooling (ASP)}
To prevent recurrent temporal dilution over the dynamic, variable-length $8$-frame video analysis windows, the framework purposefully pools the temporal dimensions utilizing an Attentive Statistics Pooling (ASP) layer. For intermediate feature vectors $h_t$, a channel-wise weighted mean $\mu$ and standard deviation $\sigma$ are generated:
\begin{equation}
    \mu = \sum_{t} \alpha_t h_t, \quad \sigma = \sqrt{\sum_{t} \alpha_t (h_t - \mu)^2}
\end{equation}
This translates an otherwise highly-variable data stream into a rigid, uniformly shaped 512-dimensional Audio-Visual embedding optimal for L2 normalization and cosine-similarity threshold bounding.

\section{Data Engineering and Hardware Constraints}

\subsection{Dataset Characteristics}
To properly gauge deepfake vulnerabilities across independent modalities, we utilized the \textbf{FakeAVCeleb v1.2} dataset. Operating globally with 21,566 video clips built from 500 unique source identities, the schema provides a massive combinatorial sandbox to simulate impersonations. Original HD videos were algorithmically normalized to $112 \times 112$ pixels to align with the MTCNN facial extractor baseline, while corresponding audio streams were downsampled to 16kHz mono formats. 

\subsection{Hardware and Deployment Environment}
To efficiently run high-volume neural matrix multiplications across the FakeAVCeleb schema, the pipeline was orchestrated leveraging Google Colab's A100 GPU clusters. To maximize processing velocity, PyTorch 2.5.1 was enforced alongside strict \texttt{cu121} Nvidia CUDA libraries. A critical engineering hurdle materialized during dependency resolution; the \texttt{torchaudio} library bindings were meticulously isolated outside of standard pip constraints to prevent catastrophic \texttt{libcudart.so} downgrades inadvertently triggered by the \texttt{facenet-pytorch} detection logic. Additionally, broken \texttt{visual\_model.model} checkpoint layers leaving 240+ untrained residual layers were rectified by standardizing instantiation structures.

\subsection{Data Scrubbing and Provisioning}
Deepfake generation repositories are notoriously plagued by corrupted MP4 headers and desynchronized audio bitrates. To guarantee a strictly unbiased 300-pair evaluation list, we constructed a dynamic constraint filter utilizing an infinite \texttt{while} loop algorithm. A given genuine-imposter validation pair was only admitted into the active cohort if both the ground truth and synthetic variations successfully bypassed overlapping MTCNN bounding-box extractions flawlessly.

\section{Exploratory Data Analysis (EDA)}
During the preliminary exploratory phase, statistical profiling clarified several foundational attributes.

\subsection{Class Balance}
Visualizations generated natively in our Jupyter environment verified extreme structural variances in the generation schemas. Real recordings (RealVideo-RealAudio) form the foundational baseline representing approximately $\sim 15\%$ of the total aggregate payload, heavily overshadowed by the combinatorial synthesized attacks.

\subsection{Univariate and Bivariate Demographic Variance}
We computed demographic parity through gender and racial categorizations explicitly tracked in the metadata. The distributions exhibit a significant slant specifically towards male-heavy subsets. When contrasted directly against generative methods, we see that Wav2Lip representations scale uniformly. However, \textit{Faceswap} is disproportionally targeted towards specific geographic subjects depending intrinsically on initial camera angles in VoxCeleb. This finding acts as a major signal that specific deepfake configurations might bypass baseline evaluations artificially efficiently based strictly on demographic data-starvation.

\subsection{Generative Methodology Profiling}
A deep analysis of the production matrices explicitly highlights that \textbf{Wav2Lip} is the exceptionally dominant generative method, encompassing over 44\% of total deepfake samples. This critically alters baseline assumptions: since Wav2Lip focuses fundamentally on achieving high-fidelity acoustic-to-lip synchronization, it inherently manipulates the exact mechanisms the RJCA cross-attention model relies on to detect temporal inconsistencies.

\section{Baseline Implementation and Preliminary Results}

For our primary baseline approach, instead of utilizing classic linear binary margin classifiers, we extracted and normalized localized spatial-temporal embeddings directly from the \textbf{Recursive Joint Cross-Attention (RJCA) Fusion} architecture's \texttt{ECAPA-TDNN} and \texttt{IResNet-18} backbones. We then computed dynamic Cosine Similarity scores over the pooled 512-dim arrays to definitively gauge if generated deepfakes inherently mimic true localized identities.

\subsection{Performance Metrics and EER Synthesis}
We systematically evaluated the baseline structure against three distinct deepfake attack permutations against a live Genuine access boundary and a Zero-Effort Imposter (Stranger) margin. Deriving insights from 299 isolated Zero-Effort Imposter comparisons allowed us to calculate a highly precise Verification Equal Error Rate (EER) threshold mapped accurately at $\mathbf{0.65}$.

Operating against this specific similarity constraint, the RJCA baseline yields the subsequent Attack Success Rates (ASR):

\begin{itemize}
    \item \textbf{S0 (Genuine Access):} 100.0\% Acceptance (Mean Score: $0.985$)
    \item \textbf{STRANGER (Zero-Effort Imposter):} 0.0\% FAR (Mean Score: $0.305$)
    \item \textbf{S1 (Fake Video, Real Audio):} 92.0\% ASR (Mean Score: $0.796$)
    \item \textbf{S2 (Real Video, Fake Audio):} 100.0\% ASR (Mean Score: $0.920$)
    \item \textbf{S3 (Fake Video, Fake Audio):} 62.0\% ASR (Mean Score: $0.608$)
\end{itemize}

While the RJCA architecture definitively repels 100\% of random stranger impersonational attacks and successfully mitigates 38\% of heavily dual-manipulated (S3) generated arrays, it remains comprehensively vulnerable to modality-specific deepfakes.

\subsection{Threshold Analysis and Critical Errors}
Analyzing the Density Probability distribution reveals critically that \textbf{S2 Attacks} (which utilize genuine visual frames maliciously paired with spoofed adversarial audio streams) are exceptionally devastating against modern multi-modal architectures, systematically achieving a 100\% system bypass rate. The generated embedding densities perfectly overlapping with the genuine verification spike fundamentally confirm that while the cross-attention layers effectively correlate audio-visual mechanisms for generic matching, they inherently lack an adversarial penalty-based discriminator mapping, ultimately empowering high-fidelity spatial-spoofs to easily puncture the baseline boundary.

\vspace{0.5cm}
\begin{figure}[htbp]
    \centering
    \includegraphics[width=\linewidth]{attack_distribution.png}
    \caption{Density Probability of RJCA verification evaluation arrays pitting deepfake simulations directly against a live Genuine enrollment. The $0.65$ dashed limitation represents the dynamically calculated EER Validation threshold limit. Note the globally catastrophic distribution intersection of the underlying S2 architecture directly against the pristine Green Genuine metric loop.}
    \label{fig:kde_plot}
\end{figure}

\section{Theoretical Attack Vectors and Mathematical Vulnerabilities}
A consequential review of our statistical Density mapping concretely established that the underlying RJCA structure natively lacks defensive logic when subjected to advanced physical or latent-space modifications.

\subsection{Black-Box Convolutional Washing}
During rigorous system evaluations, we experimented heavily with external Gaussian Convolution applications ($\sigma \times \sigma$ randomized matrix masking algorithms) operating against the deepfake imagery sets independently. This targeted black-box deployment was hypothesized to fundamentally disguise and eliminate the high-frequency pixel anomalies generally expelled natively by Wav2Lip generative matrices. 
\begin{equation}
    G(x,y) = \frac{1}{2\pi\sigma^2} e^{-\frac{x^2+y^2}{2\sigma^2}}
\end{equation}
Outcomes decisively reflected that while this spatial filtration successfully neutralized low-level generated discrepancies, it similarly obscured pivotal biological boundary constraints native to standard facial markers, causing terminal score truncation mapping closely at the $0.90$ limit threshold, thereby blocking a pure continuous overlap.

\subsection{White-Box Adversarial Exploitation (FGSM)}
To explicitly inject deepfake mathematical mapping structures higher into the $0.985$ baseline bounds, computationally aggressive adversarial logic leveraging the Fast Gradient Sign Method (FGSM) mathematically operates as a premier systemic constraint. By recursively extracting gradients from localized loss maps against visual arrays $\nabla_x J(\theta, x, y)$, an attacker programmatically identifies specific vulnerability limits natively established inside the neural network and iteratively obfuscates the underlying visual structure utilizing a targeted micro-epsilon static limit.

\section{Proposed Defense Architecture (Phase 2 Roadmap)}
Synthesized strictly derived from the expansive baseline vulnerabilities concretely mathematically established in Phase 1, it fundamentally dictates that standard temporal verification algorithms categorically cannot rely purely on binary cosine differentiation bounds measured blindly over generic attention logic modules. To counteract this flaw programmatically throughout Phase 2 progression tasks, our roadmap encompasses appending an independent \textbf{Adversarial Frame Discriminator} pipeline specifically designed to dynamically punish sub-pixel variances identified explicitly during generative synchronization mappings.

Rather than classifying gross visual similarities, an intermediary temporal objective function logic structure trained recursively upon the pooled 512-dim embedding matrices will isolate disjointed biological anomalies present naturally within synthetic constructions. Operating with an optimized \textbf{Contrastive Penalty Network}, matching similarity constraints will be inherently severely degraded if algorithmic calculations establish unnatural underlying structural variances. This theoretical defense integration strategy natively dictates decoupling rogue S2 generative deepfakes forcefully detaching their numerical thresholding and aggressively dumping the arrays firmly outside of the standardized dynamically balanced verification zone.

\begin{thebibliography}{00}
\bibitem{ref1}
G. P. Rajasekhar and J. Alam, ``Audio-Visual Person Verification based on Recursive Fusion of Joint Cross-Attention,'' \textit{Proc. IEEE 18th Int. Conf. Automatic Face and Gesture Recognition (FG)}, 2024.

\bibitem{ref2}
H. Khalid, S. Tariq, M. Kim, and S. S. Woo, ``FakeAVCeleb: A Novel Audio-Video Multimodal Deepfake Dataset,'' \textit{Proc. NeurIPS Datasets and Benchmarks Track}, 2021.

\bibitem{ref3}
X. Liu, X. Wang, M. Sahidullah, J. Patino, H. Delgado, T. Kinnunen, M. Todisco, J. Yamagishi, N. Evans, A. Nautsch, and K. A. Lee, ``ASVspoof 2021: Towards Spoofed and Deepfake Speech Detection in the Wild,'' \textit{IEEE/ACM Trans. Audio, Speech, Language Process.}, vol. 31, pp. 2507–2522, 2023, doi: 10.1109/TASLP.2023.3285283.

\bibitem{ref4}
A. Nagrani, J. S. Chung, W. Xie, and A. Zisserman, ``VoxCeleb: Large-Scale Speaker Verification in the Wild,'' \textit{Comput. Speech Lang.}, vol. 60, 2020. (VoxCeleb2: J. S. Chung et al., INTERSPEECH 2018).

\bibitem{ref5}
Z. Yan, Y. Zhang, X. Yuan, S. Lyu, and B. Wu, ``DeepfakeBench: A Comprehensive Benchmark of Deepfake Detection,'' \textit{Proc. Advances in Neural Information Processing Systems (NeurIPS) Datasets and Benchmarks Track}, vol. 36, 2023.

\bibitem{ref6}
T. Oorloff, S. Koppisetti, N. Bonettini, D. Solanki, B. Colman, Y. Yacoob, A. Shahriyari, and G. Bharaj, ``AVFF: Audio-Visual Feature Fusion for Video Deepfake Detection,'' \textit{arXiv:2406.02951}, 2024.

\bibitem{ref7}
D. A. Coccomini, G. Kordopatis-Zilos, G. Amato, R. Caldelli, F. Falchi, and S. Papadopoulos, ``MINTIME: Multi-Identity Size-Invariant Video Deepfake Detection,'' \textit{IEEE Trans. Inf. Forensics Security (TIFS)}, 2024.

\bibitem{ref8}
Q. Meng, S. Zhao, Z. Huang, and F. Zhou, ``MagFace: A Universal Representation for Face Recognition and Quality Assessment,'' \textit{Proc. IEEE/CVF Conf. Computer Vision and Pattern Recognition (CVPR)}, 2021.
\end{thebibliography}

\end{document}
